\begin{definition}[The composed physical bond transformation]%
    \label{def:peps_torus_semiregular_bond_map}
    \lean{TNLean.PEPS.torusSemiRegularBondMap}
    \leanok
    \uses{thm:peps_full_multiplicity_bond_map,
        lem:peps_bond_coordinate_transport}
    Choose positive dimensions $d_i$ and put
    $V=\bigoplus_i\mathbb C^{d_i}$ and
    $V'=\bigoplus_i(\mathbb C^{d_i}\otimes\mathbb C^{d_i})$.
    Let $F:V\otimes V\to V'\otimes V'$ be the full endpoint
    multiplicity-restoring map. For $Q:V'\to\mathbb C^G$, put
    $C_Q=Q\otimes\overline Q$. On the physical bond spaces of a torus
    with positive periods, define
    \begin{align}
        \mathcal T_Q=
            \left(\bigotimes_{e\in\mathcal B}C_Q\right)
            \left(\bigotimes_{e\in\mathcal B}F\right).
        \label{eq:peps_torus_semiregular_map}
    \end{align}
\end{definition}

\begin{theorem}[Supported overlaps under the composed bond transformation]%
    \label{thm:peps_torus_semiregular_supported_overlap}
    \lean{TNLean.PEPS.torusSemiRegularBondMap_dotProduct}
    \leanok
    \uses{def:peps_torus_semiregular_bond_map,
        thm:peps_multiplicity_bond_product_overlap,
        thm:peps_physical_product_supported_isometry,
        lem:peps_bond_coordinate_transport}
    Suppose $Q^\dagger Q=I$. For every physical vector $\psi$ and every
    $\phi\in\operatorname{ran}\bigotimes_{e\in\mathcal B}E_V$,
    \begin{align}
        \langle\mathcal T_Q\psi,\mathcal T_Q\phi\rangle
            =\langle\psi,\phi\rangle.
    \end{align}
\end{theorem}
\begin{proof}\leanok
    The product of $C_Q$ is an isometry. Remove it from the overlap,
    then apply preservation of overlaps by the product of $F$ on the
    matching-sector support.
\end{proof}

\begin{theorem}[The actual weighted torus state under a chosen Fourier intertwiner]%
    \label{thm:peps_torus_semiregular_chosen_fourier_state}
    \lean{TNLean.PEPS.torusSemiRegularBondMap_blockFourthRootWeight}
    \leanok
    \uses{def:peps_torus_semiregular_bond_map,
        thm:peps_torus_block_multiplicity_state,
        thm:peps_torus_bond_coordinate_transport}
    Let $D_i$ be matrix representations of a finite group $G$ of
    positive dimensions $d_i$. Set
    $U=\bigoplus_iD_i$, $W=\bigoplus_i d_i^{1/4}I$, and
    $L=\bigoplus_i D_i\otimes I_{d_i}$.
    Suppose the chosen matrix $Q:V'\to\mathbb C^G$ satisfies
    $U_{\mathrm{reg}}(g)=QL(g)Q^\dagger$ for every $g\in G$.
    Let $\Psi_{U,W}$ be the actual torus state of the normalized
    averaging site dressed by $W$ on its four virtual legs, and let
    $\Psi_{\mathrm{reg}}$ be the actual regular averaging-site torus state.
    After the original physical coordinates are regrouped by bonds,
    \begin{align}
        \mathcal R\Psi_{U,W}
            &\in\operatorname{ran}\bigotimes_{e\in\mathcal B}E_V,\\
        \mathcal T_Q\mathcal R\Psi_{U,W}
            &=\mathcal R\Psi_{\mathrm{reg}}.
    \end{align}
    The representation matrices and Fourier intertwiner are explicit
    hypotheses of this auxiliary statement of the physical transformation
    in \cite[Section~7]{Schuch2010PEPS}.
\end{theorem}
\begin{proof}\leanok
    The actual weighted contraction has matching-sector support, and
    the product of $F$ sends it to the actual averaging-site state of $L$.
    The product of $C_Q$ sends that state to the actual regular state
    by the matrix intertwining identity. Compose the two equalities.
\end{proof}

\begin{theorem}[An ambient physical isometry for the chosen block construction]%
    \label{thm:peps_torus_semiregular_isometric_extension}
    \lean{TNLean.PEPS.exists_isometric_torusSemiRegularBondMap}
    \leanok
    \uses{thm:peps_torus_semiregular_chosen_fourier_state,
        thm:peps_physical_product_supported_isometry,
        thm:peps_full_multiplicity_bond_isometry_extension,
        lem:peps_bond_coordinate_transport}
    Under the representation hypotheses of the chosen Fourier transformation,
    suppose also $Q^\dagger Q=I$. There is an isometry
    $T:V\otimes V\to V'\otimes V'$ such that, with
    $\mathcal U=(\bigotimes_{e\in\mathcal B}C_Q)
        (\bigotimes_{e\in\mathcal B}T)$,
    \begin{align}
        \mathcal U\mathcal R\Psi_{U,W}&=\mathcal R\Psi_{\mathrm{reg}},\\
        \langle\mathcal U\psi,\mathcal U\phi\rangle
            &=\langle\psi,\phi\rangle.
    \end{align}
    The overlap identity holds for all vectors in the original physical space.
\end{theorem}
\begin{proof}\leanok
    Extend $F$ to an isometry $T$ satisfying $TE_V=FE_V$.
    Product matrices respect multiplication, so their products agree on
    $\operatorname{ran}\bigotimes_e E_V$. The actual weighted state lies
    in this range, and hence retains the chosen Fourier state equality.
    Both physical product maps in $\mathcal U$ are isometries.
\end{proof}

\begin{theorem}[Finite-torus equivalence from a multiplicity-one semi-regular representation]%
    \label{thm:peps_torus_multiplicity_one_semiregular_equivalence}
    \lean{TNLean.PEPS.exists_isometric_multiplicityOneSemiRegular_torusState}
    \leanok
    \uses{thm:peps_regular_minimal_representation,
        thm:peps_torus_semiregular_isometric_extension}
    Let $G$ be a finite group and let the torus have positive periods.
    There are positive dimensions $d_i$ and pairwise inequivalent unitary
    irreducible representations $D_i$ such that $U=\bigoplus_iD_i$ is
    semi-regular and each irreducible sector has multiplicity one.
    Put $W=\bigoplus_i d_i^{1/4}I$ and
    $L=\bigoplus_i D_i\otimes I_{d_i}$.
    There are isometries $Q:V'\to\mathbb C^G$ and
    $T:V\otimes V\to V'\otimes V'$ satisfying
    $U_{\mathrm{reg}}(g)=QL(g)Q^\dagger$ for every $g$.
    For $\mathcal U=(\bigotimes_{e\in\mathcal B}(Q\otimes\overline Q))
        (\bigotimes_{e\in\mathcal B}T)$,
    \begin{align}
        \mathcal U\mathcal R\Psi_{U,W}&=\mathcal R\Psi_{\mathrm{reg}},\\
        \langle\mathcal U\psi,\mathcal U\phi\rangle
            &=\langle\psi,\phi\rangle.
    \end{align}
    The latter identity holds for all vectors in the original physical space.
    This is the finite-torus realization of the local physical equivalence
    in \cite[Section~7]{Schuch2010PEPS}; the representation and Fourier data
    are conclusions rather than additional hypotheses.
\end{theorem}
\begin{proof}\leanok
    Construct the orthonormal Fourier basis of the actual regular
    representation, with each irreducible repeated by its dimension.
    The associated single-copy block representation is unitary and
    semi-regular, with character multiplicity one in every sector.
    Its change to the original group basis gives $Q$ and the matrix
    intertwining relation. Apply the ambient physical isometry theorem
    to these derived data.
\end{proof}
